\section{Proof Decomposition}

The end goal is the correctness of $\mathsf{ABA.Impl}$.
We establish it via a simulation by $\mathsf{ABA.Spec}$ and transport the spec's correctness.
The simulation itself is decomposed recursively along the sub-protocol structure.

For a protocol $P$ with concrete sub-protocols $Q_1,\dots,Q_k$, the step ``$\mathsf{P.Spec}$ simulates $\mathsf{P.Impl}$'' is split into:
\begin{enumerate}
    \item (Substitution) $\mathsf{P.Impl}[\mathsf{Q_1.Spec},\dots,\mathsf{Q_k.Spec}]$ is simulated by $\mathsf{P.Impl}$, using the child simulations $\mathsf{Q_i.Spec}$ simulates $\mathsf{Q_i.Impl}$ and compositionality.
    \item (Core) $\mathsf{P.Spec}$ simulates the hybrid $\mathsf{P.Impl}[\mathsf{Q_1.Spec},\dots,\mathsf{Q_k.Spec}]$.
    \item (Transitivity) Compose (1) and (2).
\end{enumerate}
Leaves of the recursion (no sub-protocol, or only the black-box $\mathsf{AVSS.Spec}$) skip the substitution step.

\subsection{Top-level}

\begin{theorem}\label{thm:ABACorrect}
    \lean{ABA.Correct}
    %\leanok
    \uses{def:ABAConcrete, prot:SpecABA}
    $\mathsf{ABA.Concrete}$ satisfies Validity, Agreement and Almost-Sure Termination.
\end{theorem}
\begin{proof}
    \uses{thm:ABASim, thm:pfs-sound}
    Soundness of probabilistic forward simulation (Theorem~\ref{thm:pfs-sound}) applied to \ref{thm:ABASim} yields $\mathsf{ABA.Concrete} \sqsubseteq_{\mathrm{TD}} \mathsf{ABA.Spec}$, which transports the safety properties (Validity and Agreement) of $\mathsf{ABA.Spec}$ to $\mathsf{ABA.Concrete}$. Almost-Sure Termination is a probabilistic liveness property and is handled separately via the liveness-preserving strengthening of $\preceq_{\mathrm{PFS}}$.
\end{proof}

\subsection{ABA}

\begin{theorem}\label{thm:ABASim}
    \lean{ABA.Sim}
    %\leanok
    \uses{def:ABAConcrete, prot:SpecABA}
    $\mathsf{ABA.Spec}$ simulates $\mathsf{ABA.Concrete}$.
\end{theorem}
\begin{proof}
    \uses{lem:ABASubst, lem:ABACore, thm:pfs-largest-precongruence}
    Probabilistic forward simulation is in particular a pre-order (Theorem~\ref{thm:pfs-largest-precongruence}); chain \ref{lem:ABACore} with \ref{lem:ABASubst} by transitivity.
\end{proof}

\begin{lemma}\label{lem:ABASubst}
    \lean{ABA.SubstSim}
    %\leanok
    \uses{def:ABAConcrete, def:ABAHybrid}
    $\mathsf{ABA.Hybrid}$ is simulated by $\mathsf{ABA.Concrete}$.
\end{lemma}
\begin{proof}
    \uses{thm:GBCASim, thm:WCCSim, thm:pfs-precongruence}
    $\mathsf{ABA.Concrete}$ and $\mathsf{ABA.Hybrid}$ share the same skeleton $\mathsf{ABA.Impl}$ and differ only in their sub-protocol instances. Apply the pre-congruence of $\preceq_{\mathrm{PFS}}$ for parallel composition (Theorem~\ref{thm:pfs-precongruence}) to the child simulations \ref{thm:GBCASim} and \ref{thm:WCCSim}, lifted along the $\mathsf{ABA.Impl}$ context.
\end{proof}

\begin{lemma}\label{lem:ABACore}
    \lean{ABA.CoreSim}
    %\leanok
    \uses{def:ABAHybrid, prot:SpecABA}
    $\mathsf{ABA.Spec}$ simulates $\mathsf{ABA.Hybrid}$.
\end{lemma}

\subsection{WCC}

\begin{theorem}\label{thm:WCCSim}
    \lean{WCC.Sim}
    %\leanok
    \uses{def:WCCConcrete, prot:SpecWCC}
    $\mathsf{WCC.Spec}$ simulates $\mathsf{WCC.Concrete}$.
\end{theorem}
\begin{proof}
    \uses{lem:WCCSubst, lem:WCCCore, thm:pfs-largest-precongruence}
    Probabilistic forward simulation is in particular a pre-order (Theorem~\ref{thm:pfs-largest-precongruence}); chain \ref{lem:WCCCore} with \ref{lem:WCCSubst} by transitivity.
\end{proof}

\begin{lemma}\label{lem:WCCSubst}
    \lean{WCC.SubstSim}
    %\leanok
    \uses{def:WCCConcrete, def:WCCHybrid}
    $\mathsf{WCC.Hybrid}$ is simulated by $\mathsf{WCC.Concrete}$.
\end{lemma}
\begin{proof}
    \uses{thm:GatherSim, thm:RSDSim, thm:pfs-precongruence}
    $\mathsf{WCC.Concrete}$ and $\mathsf{WCC.Hybrid}$ share the same skeleton $\mathsf{WCC.Impl}$ and differ only in the instances of Gather and RSD. Apply the pre-congruence of $\preceq_{\mathrm{PFS}}$ for parallel composition (Theorem~\ref{thm:pfs-precongruence}) to the child simulations \ref{thm:GatherSim} and \ref{thm:RSDSim}, lifted along the $\mathsf{WCC.Impl}$ context.
\end{proof}

\begin{lemma}\label{lem:WCCCore}
    \lean{WCC.CoreSim}
    %\leanok
    \uses{def:WCCHybrid, prot:SpecWCC}
    $\mathsf{WCC.Spec}$ simulates $\mathsf{WCC.Hybrid}$.
\end{lemma}

\subsection{GBCA}

\begin{theorem}\label{thm:GBCASim}
    \lean{GBCA.Sim}
    %\leanok
    \uses{prot:ImplGBCA, prot:SpecGBCA}
    $\mathsf{GBCA.Spec}$ simulates $\mathsf{GBCA.Impl}$. No sub-protocol: direct proof.
\end{theorem}

\subsection{Gather}

\begin{theorem}\label{thm:GatherSim}
    \lean{Gather.Sim}
    %\leanok
    \uses{def:GatherConcrete, prot:SpecGather}
    $\mathsf{Gather.Spec}$ simulates $\mathsf{Gather.Concrete}$.
\end{theorem}
\begin{proof}
    \uses{lem:GatherSubst, lem:GatherCore, thm:pfs-largest-precongruence}
    Probabilistic forward simulation is in particular a pre-order (Theorem~\ref{thm:pfs-largest-precongruence}); chain \ref{lem:GatherCore} with \ref{lem:GatherSubst} by transitivity.
\end{proof}

\begin{lemma}\label{lem:GatherSubst}
    \lean{Gather.SubstSim}
    %\leanok
    \uses{def:GatherConcrete, def:GatherHybrid}
    $\mathsf{Gather.Hybrid}$ is simulated by $\mathsf{Gather.Concrete}$.
\end{lemma}
\begin{proof}
    \uses{thm:BRBSim, thm:pfs-precongruence}
    $\mathsf{Gather.Concrete}$ and $\mathsf{Gather.Hybrid}$ share the same skeleton $\mathsf{Gather.Impl}$ and differ only in the BRB instances. Apply the pre-congruence of $\preceq_{\mathrm{PFS}}$ for parallel composition (Theorem~\ref{thm:pfs-precongruence}) to the child simulation \ref{thm:BRBSim}, lifted along the $\mathsf{Gather.Impl}$ context (one instance per leader).
\end{proof}

\begin{lemma}\label{lem:GatherCore}
    \lean{Gather.CoreSim}
    %\leanok
    \uses{def:GatherHybrid, prot:SpecGather}
    $\mathsf{Gather.Spec}$ simulates $\mathsf{Gather.Hybrid}$.
\end{lemma}

\subsection{RSD}

\begin{theorem}\label{thm:RSDSim}
    \lean{RSD.Sim}
    %\leanok
    \uses{def:RSDConcrete, prot:SpecRSD}
    $\mathsf{RSD.Spec}$ simulates $\mathsf{RSD.Concrete}$.
\end{theorem}
\begin{proof}
    \uses{lem:RSDSubst, lem:RSDCore, thm:pfs-largest-precongruence}
    Probabilistic forward simulation is in particular a pre-order (Theorem~\ref{thm:pfs-largest-precongruence}); chain \ref{lem:RSDCore} with \ref{lem:RSDSubst} by transitivity.
\end{proof}

\begin{lemma}\label{lem:RSDSubst}
    \lean{RSD.SubstSim}
    %\leanok
    \uses{def:RSDConcrete, def:RSDHybrid}
    $\mathsf{RSD.Hybrid}$ is simulated by $\mathsf{RSD.Concrete}$. Only BRB is substituted; $\mathsf{AVSS.Spec}$ stays fixed on both sides.
\end{lemma}
\begin{proof}
    \uses{thm:BRBSim, thm:pfs-precongruence}
    $\mathsf{RSD.Concrete}$ and $\mathsf{RSD.Hybrid}$ share the same skeleton $\mathsf{RSD.Impl}$ and the same $\mathsf{AVSS.Spec}$ instances; only the BRB instances are substituted. Apply the pre-congruence of $\preceq_{\mathrm{PFS}}$ for parallel composition (Theorem~\ref{thm:pfs-precongruence}) to the child simulation \ref{thm:BRBSim}, lifted along the $\mathsf{RSD.Impl}$ context (one BRB instance per leader).
\end{proof}

\begin{lemma}\label{lem:RSDCore}
    \lean{RSD.CoreSim}
    %\leanok
    \uses{def:RSDHybrid, prot:SpecRSD}
    $\mathsf{RSD.Spec}$ simulates $\mathsf{RSD.Hybrid}$.
\end{lemma}

\subsection{BRB}

\begin{theorem}\label{thm:BRBSim}
    \lean{BRB.Sim}
    %\leanok
    \uses{prot:ImplBRB, prot:SpecBRB}
    $\mathsf{BRB.Spec}$ simulates $\mathsf{BRB.Impl}$. No sub-protocol: direct proof.
\end{theorem}
